\subsection{Differentiation under the integral sign}

We fix a locally compact topological space X and a measure \(\nu\) on X. Let E be a Banach space.

PROPOSITION 2. --- Let I be an interval of \(\mathbf{R}\) and f a mapping from X × I into E such that

\begin{enumerate}
\def\labelenumi{(\roman{enumi})}
\item
  For every \(t \in I\) , the map \(x \mapsto f(x, t)\) of X into E is \(\nu\) -integrable;
\item
  For every \(x \in X\) , the map \(t \mapsto f(x, t)\) from I into E admits a derivative, denoted \(t \mapsto f'(x, t)\) ;
\item
  There exists a positive function \(\nu\) -integrable g on X such that \(\|f'(x,t)\| \leqslant g(x)\) for all \((x,t) \in \mathrm{X} \times \mathrm{I}\) .
\end{enumerate}

Then the map F from I into E defined by

\[
\mathrm{F} (t) = \int_ {\mathrm{X}} f (x, t) d \nu (x)
\]

is differentiable in I and for every \(t \in \mathrm{I}\) , we have

\[
\mathrm{F} ^ {\prime} (t) = \int_ {\mathrm{X}} f ^ {\prime} (x, t) d \nu (x).
\]

Let \(t_0 \in \mathrm{I}\) and let \(\mathrm{J}\) be an interval of \(\mathbf{R}\) contained in \(\mathrm{I}\) which is a neighborhood of \(t_0\) in \(\mathrm{I}\) . Let \(h: \mathrm{X} \times \mathrm{J} \to \mathrm{E}\) be the map defined by \((x, t) \mapsto (f(x, t) - f(x, t_0)) / (t - t_0)\) for \((x, t) \in \mathrm{X} \times (\mathrm{J} - \{t_0\})\) and \((x, t_0) \mapsto f'(x, t_0)\) for \(x \in \mathrm{X}\) . Let \(x \in \mathrm{X}\) . We have \(\|h(x, t_0)\| \leqslant g(x)\) and, for all \(t \neq t_0\) , it follows that \(\|h(x, t)\| \leqslant g(x)\) . By FVR, I, p. 23, th. 2. By definition of the derivative, the proposition follows from Corollary 1 of INT, IV, p. 144, § 4, no. 3, applied to the mapping \(h\) .

Let us take up the notations of VAR, R1, 2.4, p. 19 concerning partial derivatives.

COROLLARY 1. --- Let U be an open subset of \(\mathbf{R}^{n}\) . Let \(k \in \mathbf{N}\) and \(f\) be a map from \(\mathrm{X} \times \mathrm{U}\) into E satisfying the following conditions:

\begin{enumerate}
\def\labelenumi{(\roman{enumi})}
\item
  For every \(t \in U\) , the map \(x \mapsto f(x, t)\) of X into E is \(\nu\) -integrable;
\item
  For every \(x \in X\) , the map \(t \mapsto f(x, t)\) from U into E is of class \(C^{k}\) , with partial derivatives denoted \(\partial^{\alpha}f(x, t)\) for all \(\alpha \in \mathbf{N}^{n}\) such that \(|\alpha| \leqslant k\) ;
\item
  There exists a function \(\nu\) -integrable \(g\) on X such that for every \(\alpha \in \mathbf{N}^n\) satisfying \(|\alpha| \leqslant k\) and for every \((x,t) \in \mathrm{X} \times \mathrm{U}\) , we have \(\|\partial^\alpha f(x,t)\| \leqslant g(x)\) .
\end{enumerate}

Then the map F from U to E defined by

\[
\mathrm{F} (t) = \int_ {\mathrm{X}} f (x, t) d \nu (x)
\]

is of class \(C^{k}\) in U and for all \(\alpha\in \mathbf{N}^{n}\) satisfying \(|\alpha|\leqslant k\) and every \(t\in U\) , we have

\[
\partial^ {\alpha} \mathrm{F} (t) = \int_ {\mathrm{X}} \partial^ {\alpha} f (x, t) d \nu (x).
\]

This follows from the proposition by induction on k.

COROLLARY 2. --- Let k be a natural number. Let \(f \in \mathcal{L}^1(\mathbf{R}^n, \mu)\) and \(g \in \mathcal{C}^k(\mathbf{R}^n)\) . Assume that for every \(\alpha \in \mathbf{N}^n\) such that \(|\alpha| \leqslant k\) , the function \(\partial^\alpha g\) is bounded. Then the convolution product \(f * g\) belongs to \(\mathcal{C}^k(\mathbf{R}^n)\) and for every \(\alpha\) such that \(|\alpha| \leqslant k\) , we have \(\partial^\alpha(f * g) = f * \partial^\alpha g\) .

We can apply Corollary 1 to the space \(X = \mathbf{R}^{n}\) with Lebesgue measure and to the mapping h defined by \((x, t) \mapsto f(x)g(t - x)\) from

\(\mathbf{R}^{n} \times \mathbf{R}^{n}\) into \(\mathbf{C}\); indeed, for every \(\alpha \in \mathbf{N}^{n}\) such that \(|\alpha| \leqslant k\) , we have the inequality

\[
| \partial^ {\alpha} h (x, t) | \leqslant \Bigl (\sup _ {| \beta | \leqslant k} \sup _ {y \in \mathbf {R} ^ {n}} | \partial^ {\beta} g (y) | \Bigr) | f (x) |
\]

whose right-hand side is an integrable function on \(R^{n}\) .

